\subsubsection{问题四具体分析}

问题四要求同时优化任务时序与储能充放电，并估计"仅任务侧""仅储能侧""双侧联合"三者的边界与交互。本文采用四格分解：固定同一可行任务族与同一能源边界，逐步开关两类灵活性，得到 F00（固定+无储能）、F10（任务+无储能）、F01（固定+储能）、F11（任务+储能）四个策略价值，并以 $I=F_{10}+F_{01}-F_{11}-F_{00}$ 度量交互（$I>0$ 表示互补、$I<0$ 表示替代）。随后通过新能源规模缩减 $\rho\in\{1.0,0.9,\dots,0.4\}$、短时新能源缺口与电价机制扰动三组情景检验鲁棒性。

需要说明：任务侧为启发式可行策略（无有效全问题下界），储能侧为子问题最优，故 $I$ 为可行策略的点估计，不报告认证区间。四格物质保证：共用同一可行任务族与同一能源边界；F01 在 Q4 公共基线（本地调度 aiN）上重解储能，与 F00 同负荷体系，而非复用问题三的附件负荷。为保证可行域嵌套（$\Omega_{00}\subseteq\Omega_{10}\subseteq\Omega_{11}$、$\Omega_{01}\subseteq\Omega_{11}$），F10 保留"不迁移"incumbent（$F_{10}\le F_{00}$），F11 在"两阶段(迁移+储能)、F01、F10"三者取当前最好（$F_{11}\le\min\{F_{01},F_{10}\}$）；存储感知邻域回退搜索\cite{pisinger2010lns}（storage-aware rollback neighborhood search）进一步从两阶段方案批量回退迁移任务、重解储能子问题加以验证。
